% ---------- ¶2  "to learn about X, we examine Y" ----------
To examine whether unraveling describes a market where quality is verifiable but silence is
common, we study the French residential housing market around the July~2021 enforcement of
Energy Performance Certificate (EPC) disclosure. Under a 2018 reform that took effect on
July~1, 2021, a property's \emph{diagnostic de performance \'energ\'etique} (DPE) became
legally opposable.\footnote{Article 179 III of \emph{loi n\textdegree{}2018-1021 du 23 novembre
2018} (loi \'ELAN) amended Articles L271-4 and L134-3-1 of the \emph{Code de la construction et
de l'habitation}, with opposability effective July~1, 2021; the accompanying methodology
overhaul (replacing the earlier invoice-based calculation with the current
physical-characteristics-based method) was set by the \emph{Arr\^et\'e du 31 mars 2021}. The
\emph{loi Climat et R\'esilience} (\emph{loi n\textdegree{}2021-1104 du 22 ao\^ut 2021}) is a
separate, later reform; we invoke it elsewhere in the paper for the F/G rental-ban calendar it
introduced, not for DPE opposability.} A seller who advertises an energy score is now bound by
it, and buyers acquired legal recourse against misstatement. The reform changed what a disclosed
score means without changing the underlying housing stock or the legal consequences of remaining
silent. It therefore isolates the demand-side value of the signal, and it does so in a market
where the certificate is issued and archived by a national agency rather than by the seller or
the listing platform.
